\documentclass[14pt,dvipdfmx]{jsarticle}
\input{lecture-preamble.tex}

\begin{document}

\lectureheader{数学Ⅱ}{5}{指数関数と対数関数}{2}{対数関数}{3-BC}{対数の性質}
  {163,164}{対数の性質と底の変換を使い, 対数を計算できるようになろう}

% 同じ節の続き（3-A までの番号を引き継ぐ）
\setcounter{definition}{1}
\setcounter{theorem}{2}

% ============ 1ページ目（表・左） ============
\vspace{-1.8zw}%
\heading{〜前時の復習〜}
\begin{reviewbox}{対数の定義と基本}
  \begin{enumerate}[label=\textbf{\arabic*}, leftmargin=1.5zw, itemsep=0.02em]
    \item $M = a^p \iff p = \fitblankbf{\log_{a} M}$
    \item $a$ を\fitblankbf{\text{底}}, $M$ を\fitblankbf{\text{真数}}という.

    真数 $M$ は常に\fitblankbf{\text{正の数}}である.
    \item $\log_{a} a^{p} = \fitblankbf{p}$
  \end{enumerate}
\end{reviewbox}

$\fitblankbf{1} = a^{0}$, $\fitblankbf{a} = a^{1}$ であることから, 次が成り立つ.
\begin{theorembox}{特別な値}
  {\centering
  $\log_{a} 1 = \fitblankbf{0}$,\quad $\log_{a} a = \fitblankbf{1}$
  \par}
\end{theorembox}

\begin{theorembox}{対数の性質}
  $M > 0$, $N > 0$ で, $k$ は実数とする.
  \par\vspace{0.2em}
  \begin{enumerate}[label=\textbf{\arabic*}, leftmargin=1.5zw, itemsep=0.40em, topsep=0.15em]
    \item $\log_{a} MN = \fitblankbf{\log_{a} M + \log_{a} N}$
    \item $\log_{a} \dfrac{M}{N} = \fitblankbf{\log_{a} M - \log_{a} N}$
    \item $\log_{a} M^{k} = \fitblankbf{k \log_{a} M}$
  \end{enumerate}
\end{theorembox}

\vspace{0.08em}
\textbf{注意}\quad 2 において, とくに $M=1$ のときは\quad
{\boldmath $\log_{a} \dfrac{1}{N} = -\log_{a} N$}

\vfill
\nextpage

% ============ 2ページ目（表・右） ============
\vspace{-1.5zw}%
\textbf{【1 の証明】}
\begin{answerbox}{1}[8]
  $\log_{a} M = p$, $\log_{a} N = q$ とすると\\
  $MN = a^{p} \times a^{q} = a^{p+q}$\\
  したがって
  \boxanswer{$\log_{a} M + \log_{a} N$}
\end{answerbox}

\begin{practiceprob}{16}
  上の証明にならって、性質2、3が成り立つことを証明せよ。
\end{practiceprob}

\begin{answerbox}{2}[8]
  $\log_{a} M = p$, $\log_{a} N = q$ とすると\\
  $M = a^{p}$, $N = a^{q}$ より $\dfrac{M}{N} = a^{p-q}$
  \boxanswer{$\log_{a} M - \log_{a} N$}
\end{answerbox}
\begin{answerbox}{3}[8]
  $\log_{a} M = p$ とすると $M = a^{p}$\\
  両辺を $k$ 乗すると $M^{k} = a^{pk}$
  \boxanswer{$k \log_{a} M$}
\end{answerbox}

\begin{exampleprob}{10}
  次の式を計算せよ。
  \begin{tasks}[label=(\arabic*), label-width=1.6zw, item-indent=2zw, after-item-skip=0.1em](2)
    \task $\log_{10} 2 + \log_{10} 5$
    \task $\log_{2} 24 - \log_{2} 3$
    \task $2\log_{3} 4 + \log_{3} 5 - \log_{3} 8$
  \end{tasks}
\end{exampleprob}

\begin{answerbox}{(1)}[5]
  $\log_{10}(2 \times 5)$
  \boxanswer{$1$}
\end{answerbox}
\begin{answerbox}{(2)}[5]
  $\log_{2} \dfrac{24}{3}$
  \boxanswer{$3$}
\end{answerbox}
\begin{answerbox}{(3)}[5]
  $\log_{3} \dfrac{4^{2} \times 5}{8}$
  \boxanswer{$\log_{3} 10$}
\end{answerbox}

\vfill
\nextpage

% ============ 3ページ目（裏・左） ============
\vspace{-1.5zw}%
\begin{practiceprob}{17}
  次の式を計算せよ。
  \begin{tasks}[label=(\arabic*), label-width=1.6zw, item-indent=2zw, after-item-skip=0.12em](2)
    \task $\log_{6} 3 + \log_{6} 12$
    \task $\log_{3} 2 - \log_{3} 18$
    \task $\log_{3} 4 + \log_{3} 18 - 3\log_{3} 2$
    \task $\log_{5} 12 - \log_{5} 3 - 2\log_{5} 10$
  \end{tasks}
\end{practiceprob}

\begin{answerbox}{(1)}[6]
  $\log_{6}(3 \times 12) = \log_{6} 36$
  \boxanswer{$2$}
\end{answerbox}
\begin{answerbox}{(2)}[6.5]
  $\log_{3} \dfrac{2}{18} = \log_{3} \dfrac{1}{9}$
  \boxanswer{$-2$}
\end{answerbox}
\begin{answerbox}{(3)}[7]
  $\log_{3} \dfrac{4 \times 18}{2^{3}} = \log_{3} 9$
  \boxanswer{$2$}
\end{answerbox}
\begin{answerbox}{(4)}[7]
  $\log_{5} \dfrac{12}{3 \times 10^{2}} = \log_{5} \dfrac{1}{25}$
  \boxanswer{$-2$}
\end{answerbox}

\vfill
\nextpage

% ============ 4ページ目（裏・右） ============
\vspace{-1.8zw}%
\heading{底の変換}
$a$ を底とする対数 $\log_{a} b$ を, $c$ を底とする対数で表してみよう.
\par\vspace{0.12em}
$\log_{a} b = p$ とすると $b = a^{p}$.
$c$ を底とする両辺の対数をとると $\log_{c} b = p \log_{c} a$.\\
よって $p = \fitblankbox{\dfrac{\log_{c} b}{\log_{c} a}}$.

\begin{theorembox}{底の変換}
  $a$, $b$, $c$ は正の数で, $a \neq 1$, $c \neq 1$ とするとき
  \par\vspace{0.15em}
  {\centering
  $\log_{a} b = \fitblankbox{\dfrac{\log_{c} b}{\log_{c} a}}$
  \par}
\end{theorembox}

\begin{exampleprob}{11}
  次の式を計算せよ。
  \begin{tasks}[label=(\arabic*), label-width=1.6zw, item-indent=2zw, after-item-skip=0.08em](2)
    \task $\log_{8} 16$
    \task $\log_{2} 3 \cdot \log_{3} 8$
  \end{tasks}
\end{exampleprob}

\noindent
\begin{minipage}[t]{0.49\linewidth}
\begin{answerbox}{(1)}[5.8]
  $\dfrac{\log_{2} 2^{4}}{\log_{2} 2^{3}}$
  \boxanswer{$\textstyle\dfrac{4}{3}$}
\end{answerbox}
\end{minipage}\hfill
\begin{minipage}[t]{0.49\linewidth}
\begin{answerbox}{(2)}[5.8]
  $\log_{2} 3 \cdot \dfrac{\log_{2} 8}{\log_{2} 3}$
  \boxanswer{$3$}
\end{answerbox}
\end{minipage}

\begin{practiceprob}{18}
  次の式を簡単にせよ。
  \begin{tasks}[label=(\arabic*), label-width=1.6zw, item-indent=2zw, after-item-skip=0.08em](3)
    \task $\log_{4} 8$
    \task $\log_{9} 3$
    \task $\log_{2} 3 \cdot \log_{3} 27$
  \end{tasks}
\end{practiceprob}

\noindent
\begin{minipage}[t]{0.32\linewidth}
\begin{answerbox}{(1)}[5.5]
  $\dfrac{\log_{2} 8}{\log_{2} 4}$
  \boxanswer{$\textstyle\dfrac{3}{2}$}
\end{answerbox}
\end{minipage}\hfill
\begin{minipage}[t]{0.32\linewidth}
\begin{answerbox}{(2)}[5.5]
  $\dfrac{\log_{3} 3}{\log_{3} 9}$
  \boxanswer{$\textstyle\dfrac{1}{2}$}
\end{answerbox}
\end{minipage}\hfill
\begin{minipage}[t]{0.32\linewidth}
\begin{answerbox}{(3)}[5.5]
  $\log_{2} 27$
  \boxanswer{$3$}
\end{answerbox}
\end{minipage}

\begin{summarybox}
  \begin{enumerate}[label=\textbf{\arabic*}, leftmargin=1.5zw, itemsep=0.02em]
    \item $\log_{a} MN = \fitblankbf{\log_{a} M + \log_{a} N}$
    \item $\log_{a} \dfrac{M}{N} = \fitblankbf{\log_{a} M - \log_{a} N}$,\quad
          $\log_{a} M^{k} = \fitblankbf{k \log_{a} M}$
    \item $\log_{a} b = \fitblankbox{\dfrac{\log_{c} b}{\log_{c} a}}$
  \end{enumerate}
\end{summarybox}

\vfill
\nexttime{対数関数のグラフをかき, 増加・減少を調べてみよう!}

\end{document}
